\begin{afterword}
\summaryhead

A short story. Jeffreyssai, a teacher of the beisutsukai, asks five students how the scientists of the ``Eld'' era, our own, failed for over thirty years to see what quantum physics implied, when not even Einstein, Schrödinger or von Neumann saw it. The students try religion, too few physicists, administrative burdens and a lack of practical need. Brennan argues that the Manhattan Project had a use in view and quantum theory did not.

The teacher's own answer is that Eld science thought it acceptable to take thirty years, because its process aimed at the truth eventually. He drills Brennan with rapid questions, punishes an answer given because it fits the pattern, and argues from the number of minutes in a month that a major problem might be solved in weeks. Put on the spot for five minutes, Brennan adds that Eld scientists were trained to use known science, not to resolve a major confusion, which they met once a lifetime and faced as amateurs. The teacher accepts this and sets the class to find one more flaw by the next day.

\responsehead

The story is fiction and is judged for what it shows. It stages the thesis this sequence goes on to argue, that science is too slow and a trained mind can be faster. The story never names what the founders missed, but its talk of a ``universal law'' points to the verdict of ``Many Worlds, One Best Guess,'' posted the day before: that the thirty years were a failure to see many-worlds. The story treats that verdict as settled history. The notes on that post find it disputed, so the story draws its lesson from a failure that is itself in question.

Its best scene is the drill. After three questions whose answer is ``less than a second,'' the fourth has the same form and a different answer, and the student who answers by pattern is shouted down. Brennan then refuses to answer before he has thought, and is praised for it. As in ``Initiation Ceremony,'' the test is whether a student gives the answer he has worked out or the one the teacher seems to expect. The small facts also hold. The arithmetic of minutes and insights is correct, a grant agency's six months matches the US National Science Foundation's goal at the time, and Schrödinger's retreat with a mistress matches the usual account.

The story's history of the founders is weaker. Schrödinger came close to what the teacher says no one saw: in 1952 he said that the alternatives in his equation ``all really happen simultaneously.'' Brennan's claim that resolving a major confusion was done ``once per lifetime if you were lucky'' sits uneasily with Einstein, who had done it with special and general relativity, and with the quarter century of the old quantum theory, a period of confusion in which the founders worked. Hiriwa's ``no probability theory'' overlooks Jeffreys's Bayesian \textsc{Theory of Probability} of 1939.

The one-month argument shows how many minutes a month holds. How many insights a major problem needs is the teacher's assertion, and ``it has been done'' is given no example beyond Schrödinger's retreat.

In short: a lively classroom scene whose drill against answers given by pattern is well made and whose small facts check, built on the contested many-worlds verdict as settled history, with an account of the founders as amateurs at confusion that sits uneasily with their record.
\end{afterword}
